
\section{11. Conclusion}

\subsection{11.1 Summary}

We presented \textbf{HALO-ReGLA}, a hybrid attention module that runs dense 
softmax and rectified linear streams in parallel, combined by per-head 
learnable gates. The gate is initialized to favor dense behavior and shifts 
toward linear during training.

\subsection{11.2 Results Recap}

\begin{table}[H]
\centering
\small
\begin{tabular}{llll}
\toprule
\textbf{Model} & \textbf{Acc $\Delta$} & \textbf{Mean gate} & \textbf{VRAM $\Delta$} \\
\midrule
DistilBERT & -1.2\% & 0.72 & --15\% \\
RoBERTa & -2.1\% & 0.68 & --8\% \\
GPT-2 & -1.9\% & 0.65 & --5\% \\
\bottomrule
\end{tabular}
\end{table}

\subsection{11.3 Key Takeaways}

1. \textbf{Learned gates are bimodal}: dense vs linear heads.
2. \textbf{Layer depth correlates with gate.}
3. \textbf{Learned fusion beats fixed fusion.}
4. \textbf{ELU+1 is optimal kernel.}
5. \textbf{No speedup without hard routing.}

\subsection{11.4 Future Work}

\subsubsection{Short-Term}
\begin{itemize}
\item Hard routing via thresholded gates
\item Chunked cumsum for long sequences
\item Learned kernels
\item Long-context evaluation
\end{itemize}

\subsubsection{Long-Term}
\begin{itemize}
\item Theoretical error bounds
\item Integration with GQA / KV cache
\item Cross-task gate transfer
\end{itemize}

\subsection{11.5 Closing Remarks}

HALO-ReGLA provides an interpretable hybrid attention mechanism. We release 
code and checkpoints.

\subsection{11.6 Acknowledgments}

We thank {{ACKNOWLEDGMENTS}}.

\subsection{11.7 Author Contributions}

\begin{itemize}
\item \textbf{{{AUTHOR\_1}}} designed HALO-ReGLA, implemented the code, wrote the paper.
\item \textbf{{{AUTHOR\_2}}} provided feedback, edited.
\end{itemize}

\subsection{11.8 Conflict of Interest}

None.